\subsubsection{Li Chao Tree (long long)}
\begin{lstlisting}[style=mystyle, language=C++]
// TC: O(log X) per query/insert
// usa: mantener un lower envelope de lineas y consultar en cualquier x
#include <bits/stdc++.h>
using namespace std;

const long long LINF = (1LL<<60);

struct Line{
	long long m, b; // y = m*x + b
	Line(long long m=0, long long b=LINF): m(m), b(b){}
	long long get(long long x) const { return m*x + b; }
};

struct LiChao{
	struct Node{
		Line line;
		Node *l, *r;
		Node(const Line& line): line(line), l(nullptr), r(nullptr){}
	};
	long long X_LO, X_HI;
	Node* root;
	LiChao(long long X_LO, long long X_HI): X_LO(X_LO), X_HI(X_HI){
		root = nullptr;
	}
	void addLine(Line newLine){ addLine(root, X_LO, X_HI, newLine); }
	void addLine(Node*& node, long long lo, long long hi, Line newLine){
		if(!node){
			node = new Node(newLine);
			return;
		}
		long long mid = (lo + hi) >> 1;
		bool left = newLine.get(lo) < node->line.get(lo);
		bool middle = newLine.get(mid) < node->line.get(mid);
		if(middle){
			swap(node->line, newLine);
		}
		if(hi - lo == 0) return;
		if(left != middle){
			addLine(node->l, lo, mid, newLine);
		}else{
			addLine(node->r, mid + 1, hi, newLine);
		}
	}
	long long query(long long x){ return query(root, X_LO, X_HI, x); }
	long long query(Node* node, long long lo, long long hi, long long x){
		if(!node) return LINF;
		long long res = node->line.get(x);
		if(lo == hi) return res;
		long long mid = (lo + hi) >> 1;
		if(x <= mid) return min(res, query(node->l, lo, mid, x));
		else return min(res, query(node->r, mid + 1, hi, x));
	}
};

int main(){
	LiChao lc(-(long long)1e9, (long long)1e9);
	lc.addLine(Line(1, 0));  // y = x
	lc.addLine(Line(-1, 5)); // y = -x + 5
	cout << "min at x=2: " << lc.query(2) << "\n";   // -2+5 = 3
	cout << "min at x=10: " << lc.query(10) << "\n"; // 10 vs -5 -> -5
	return 0;
}
\end{lstlisting}
